\documentclass[11pt,a4paper]{article}
\usepackage[margin=1in]{geometry}
\usepackage{hyperref}
\usepackage{enumitem}
\usepackage{graphicx}
\usepackage{xcolor}

% -----------------------------------------------------
% bvraghav-tiet-ucs503p-article.sty
% -----------------------------------------------------

% Variable: Lab Instructor
\makeatletter \newcommand{\BvrSetLabInstructor}[1]{%
  \gdef\Bvr@LabInstructor{#1}}
% default
\newcommand{\Bvr@LabInstructor}{Dr. Raghav %
  B. Venkataramaiyer}
% public accessor
\newcommand{\BvrLabInstructorValue}{\Bvr@LabInstructor}
\makeatother

% Add subtitle
\usepackage{titling}
% -- Set title bold
\pretitle{\begin{center}\bfseries\fontsize{18bp}{18bp}\selectfont}
\makeatletter
\newcommand{\BvrSubtitle}[1]{%
  \posttitle{%
    \par\end{center}
    \begin{center}\large#1\end{center}
    \vskip 0.5em}%
}
\makeatother

% Hints in the section title(s)
\newcommand{\BvrHint}[1]{%
  {\normalfont\normalsize\itshape\hfill #1}}

% -----------------------------------------------------

\title{ClassLink: A BLE-Based Classroom Attendance System}

\BvrSetLabInstructor{Dr. Raghav B. Venkataramaiyer}

\BvrSubtitle{%
  Project Proposal for UCS503P  \\
  Submitted to: \BvrLabInstructorValue \\ \bfseries
  Thapar Institute of Engineering and Technology%
}

\author{
  Vemula Manvi Smaran, CSED%
  \thanks{Roll No: \texttt{1024160111}}
  \and
  Vrishank Nehru, CSED%
  \thanks{Roll No: \texttt{1024160109}}
  \and
  Mahi Singh, CSED%
  \thanks{Roll No: \texttt{1024160102}}
  \and
  Khagendra Saini, CSED%
  \thanks{Roll No: \texttt{1024160123}}
}

\date{\today}

\begin{document}
\maketitle

\section{Higher-order goal%
  \BvrHint{\ldots secondary framing}}
This project aims to reduce administrative overhead and time lost to
manual attendance-taking in classrooms, freeing instructional time and
producing a reliable, auditable attendance record. The broader framing is
\emph{efficient, low-friction classroom administration}; the immediate
engineering objective is to translate that into a measurable, deployable
mobile/backend system.

\section{Time-to-value %
\BvrHint{\ldots primary heuristic}}
\begin{itemize}[leftmargin=0.7cm]
    \item We deliver meaningful increments quickly: an offline
    proof-of-concept validating the core Bluetooth Low Energy (BLE)
    proximity mechanism was built first, before any backend or auth work.
    \item We prioritize fast feedback over perfection by working in
    vertical slices (auth $\rightarrow$ classrooms $\rightarrow$ device
    binding $\rightarrow$ BLE $\rightarrow$ attendance transaction
    $\rightarrow$ teacher operations), keeping the app runnable after every
    phase.
    \item Success in the first iteration is defined narrowly: prove that a
    teacher phone can advertise a session over BLE and a student phone can
    detect it by proximity and check in, with a working CI pipeline
    running analyzers and tests on every change.
\end{itemize}

\section{Problem Statement}
Manual, roll-call-based classroom attendance in a college setting is slow,
easy to proxy, and produces attendance data that is hard to audit or
export. Common pain points include:
\begin{itemize}[leftmargin=0.7cm]
    \item \textbf{Time cost:} calling out or circulating a sheet consumes
    several minutes of every lecture.
    \item \textbf{Proxy attendance:} a physical sheet or verbal roll call
    cannot verify that the responding student is actually present.
    \item \textbf{No live visibility:} the teacher has no real-time
    roster of who has checked in during the session.
    \item \textbf{Poor exportability:} paper or ad-hoc digital records are
    difficult to aggregate into per-course, per-student attendance
    percentages for reporting.
\end{itemize}

\textbf{Why this matters now (contextual relevance):} the problem is
directly experienced by our own project team as students, in our own
classrooms, making it both immediately relevant and immediately
testable with real hardware and real class sessions.

\section{Proposed Solution %
\BvrHint{\ldots Software Proposition}}
\subsection{Overview}
Build \textbf{ClassLink}, an Android-first Flutter mobile application
backed by a FastAPI service and a Supabase/PostgreSQL database, with the
following components:
\begin{itemize}[leftmargin=0.7cm]
    \item \textbf{Teacher app role:} create classrooms, share a join code,
    start/stop a BLE attendance session, and view a live roster.
    \item \textbf{Student app role:} join a classroom once, then scan for
    an active session's BLE beacon and submit a signed attendance claim.
    \item \textbf{Server-authoritative backend:} validates enrollment,
    session state, token freshness, device binding, duplicate claims, and
    proximity evidence before accepting an attendance record.
    \item \textbf{Audit trail:} every security-sensitive action (device
    binding, manual correction, export) is logged as an audit event.
\end{itemize}

\subsection{Core workflow}
\begin{enumerate}[leftmargin=0.7cm]
    \item Teacher starts an attendance session; the phone begins
    advertising a short-lived, rotating BLE token identifying the
    session.
    \item Enrolled student phones nearby scan for the token, collect
    several RSSI (signal strength) samples, and apply a configured
    proximity rule.
    \item On a locally authenticated device, the student's phone signs and
    submits an attendance claim to the backend.
    \item The backend validates the claim end-to-end and returns a
    server-authoritative result; the teacher's live roster updates
    accordingly.
\end{enumerate}

\subsection{Operational constraints for an educational, single-classroom
setting}
\begin{itemize}[leftmargin=0.7cm]
    \item Target Android phones as the primary, tested platform; keep
    abstractions portable without claiming untested iOS support.
    \item No Bluetooth mesh or student-to-student relay in the first
    release --- strictly teacher-advertises, student-scans.
    \item Keep the BLE advertisement payload small (at or below 20 bytes)
    and free of any student/teacher identifiers.
\end{itemize}

\section{Solution Approach %
\BvrHint{\ldots Engineering focus}}
This is an engineering course project, so we focus on system correctness,
security boundaries, and reproducible evaluation rather than novel
algorithm research.

\subsection{Proximity assistance %
\BvrHint{\ldots non-research}}
Proximity decisions are heuristic, not exact positioning:
\begin{itemize}[leftmargin=0.7cm]
    \item Treat RSSI as noisy evidence, never as proof of exact distance.
    \item Use the median of multiple observations within a sampling
    window rather than deciding on a single packet.
    \item Store the RSSI threshold and sample policy per session so
    experiments remain reproducible; do not present preset defaults as
    scientific findings.
\end{itemize}

\subsection{Mobile + backend architecture %
\BvrHint{\ldots BLE-based solution}}
A three-tier architecture:
\begin{itemize}[leftmargin=0.7cm]
    \item \textbf{Mobile client (Flutter/Riverpod):} teacher
    advertiser/student scanner roles, local biometric/device
    authentication, offline-safe state handling.
    \item \textbf{Backend API (FastAPI):} authentication verification,
    session/token lifecycle, claim validation pipeline, roster and export
    endpoints.
    \item \textbf{Database (Supabase/PostgreSQL):} classrooms,
    enrollments, sessions, tokens, attendance records, and audit events,
    protected by row-level security policies.
\end{itemize}

\subsection{CI/CD and fast delivery enabler}
CI/CD will be used to:
\begin{itemize}[leftmargin=0.7cm]
    \item Run \texttt{flutter analyze}, \texttt{dart format}, and Flutter
    unit/widget tests on every change.
    \item Run Ruff, MyPy, and Pytest for the backend, plus a SQL
    migration smoke test where feasible.
    \item Check that no \texttt{.env} secret files are committed.
\end{itemize}

\section{Evaluation Criterion %
\BvrHint{\ldots measurable and attributable}}
We define success using measurable metrics tied to the core workflow,
gathered from executed tests on real Android hardware --- never
pre-filled or assumed values.

\subsection{Primary evaluation metric}
\textbf{Detection success rate:} the fraction of labelled-position scan
attempts (front, middle, back, door, corridor, adjacent room) that
correctly accept or reject attendance according to the configured
proximity policy.

\subsection{Secondary evaluation metrics}
\begin{itemize}[leftmargin=0.7cm]
    \item \textbf{Discovery latency:} time from teacher advertising start
    to first valid student packet.
    \item \textbf{Attendance completion latency:} time from scan start to
    server-authoritative result.
    \item \textbf{False accept / false reject rate:} incorrect
    accept/reject decisions at out-of-room and in-room test positions,
    respectively.
    \item \textbf{API latency:} p50/p95 latency for the challenge and
    claim endpoints.
    \item \textbf{Duplicate and expiry enforcement:} correct rejection
    rate for duplicate claims and expired tokens.
\end{itemize}

\subsection{Pilot validation plan %
\BvrHint{\ldots high-level}}
\begin{itemize}[leftmargin=0.7cm]
    \item Run scenario tests with one teacher phone (advertiser) and at
    least two student phones (scanners) of different models, in an actual
    classroom.
    \item Record raw, timestamped CSV measurements for every scenario
    under \texttt{docs/evaluation/}, including phone model, Android
    version, room label, and position label, before deriving any summary
    tables or figures.
\end{itemize}

\section{Scalability %
\BvrHint{\ldots with foresight}}
The proposition should scale in both theoretical and practical senses.

\begin{enumerate}[leftmargin=0.7cm]
    \item \textbf{Scalable (theoretically):} the system should support
    growth in number of classrooms, sessions, and concurrent scanning
    students by:
    \begin{itemize}[leftmargin=0.7cm]
        \item keeping BLE payloads compact and session-scoped,
        \item indexing session/token lookups by alias and status,
        \item making attendance writes idempotent under retry.
    \end{itemize}

    \item \textbf{Operationally deployable (practically):} the system
    should run on common infrastructure:
    \begin{itemize}[leftmargin=0.7cm]
        \item a managed PostgreSQL instance (Supabase),
        \item a standard FastAPI/Uvicorn deployment,
        \item a CI pipeline gating merges on analyzer, lint, and test
        results.
    \end{itemize}
\end{enumerate}

\section{Engine availability heuristic}
A mature set of ``engines'' (frameworks/services) is available to
implement this as a mobile + backend project:
\begin{itemize}[leftmargin=0.7cm]
    \item Flutter/Dart for a single Android-first mobile codebase.
    \item A maintained BLE package supporting both central (scan) and
    peripheral (advertise) roles.
    \item FastAPI for a typed, testable REST backend.
    \item Supabase for managed PostgreSQL, auth, and row-level security.
    \item Pytest/Flutter test tooling for unit, widget, and integration
    tests.
\end{itemize}
We use these mature components to focus on workflow correctness, security
boundaries, and measurement rather than inventing new infrastructure.

\section{Project scope and deliverables %
\BvrHint{\ldots iteration-friendly}}
\subsection{Initial deliverable %
\BvrHint{\ldots first iteration}}
\begin{itemize}[leftmargin=0.7cm]
    \item Offline BLE proof-of-concept: teacher advertise, student scan
    and check-in by proximity, CSV export (\textbf{complete}).
    \item Supabase schema, authentication, roles, and RLS isolation tests.
    \item Classroom creation, join-code enrollment, teacher/student UI.
    \item CI: Flutter analyzer/tests and backend Ruff/MyPy/Pytest.
\end{itemize}

\subsection{Subsequent deliverables}
\begin{itemize}[leftmargin=0.7cm]
    \item Device-key binding (Android Keystore) and signed attendance
    claims.
    \item Rotating-token BLE sessions, server-authoritative claim
    validation, and live teacher roster.
    \item Manual correction workflow with mandatory reason and audit
    trail.
    \item CSV export and real-device evaluation dataset under
    \texttt{docs/evaluation/}.
\end{itemize}

\section{Risks and mitigations}
\begin{itemize}[leftmargin=0.7cm]
    \item \textbf{BLE peripheral-mode reliability on Android:} test only
    on physical devices; emulators do not reliably support GATT-server/
    peripheral advertising.
    \item \textbf{Proxy/spoofing risk:} bind attendance claims to a
    device-generated signing key and require local biometric
    authentication before signing, rather than trusting a bare BLE
    detection.
    \item \textbf{Evaluation measurement ambiguity:} fix the metric
    definitions and CSV schema (Section 7) before running any device
    experiments.
    \item \textbf{Secret leakage:} keep the session-token HMAC secret and
    Supabase service-role key server-side only; commit only
    \texttt{.env.example} files.
\end{itemize}

\section{Summary}
This project proposes ClassLink, a BLE-based classroom attendance system
that replaces manual roll call with a device-bound, server-authoritative
check-in flow. It meets the course criteria by emphasizing time-to-value
through an already-working offline proof-of-concept, implementing CI to
support frequent safe releases, and using measurable, real-device
evaluation criteria rather than assumed results. It remains focused on
engineering correctness and security boundaries rather than research-
driven novelty.

\end{document}
